\documentclass[11pt]{article}

\usepackage[margin=1in]{geometry}
\usepackage[T1]{fontenc}
\usepackage{lmodern}
\usepackage{microtype}
\usepackage{graphicx}
\usepackage{booktabs}
\usepackage{enumitem}
\usepackage{hyperref}
\usepackage{tikz}
\usepackage{caption}
\usepackage{subcaption}

\newcommand{\ExecutionRoot}{\textsc{ExecutionRoot}}
\newcommand{\BundleRoot}{\textsc{BundleRoot}}

\title{Cogitator: Deterministic Evaluation Harness with Dual Cryptographic Roots}
\author{Cogitator Authors}
\date{Version v1.0}

\begin{document}
\maketitle

\begin{abstract}
Cogitator is a deterministic evaluation harness that emits a self-contained run bundle
and two distinct cryptographic commitments. \ExecutionRoot commits to witnessed metadata
and a deterministically ordered semantic event stream, enabling semantic drift detection.
\BundleRoot commits to the entire artifact bundle for integrity checking and CI gating.
This paper defines the two roots precisely, separates witnessed metadata from provenance
metadata, and specifies verification procedures for both non-agent and agent modes.
\end{abstract}

\section{Introduction}
Cogitator targets repeatable, inspectable evaluations. Its design goal is simple:
for a fixed binary and fixed inputs, produce byte-stable artifacts and deterministic
cryptographic roots. The system treats external tool calls (including LLM inference)
as record/replay boundaries and relies on canonical JSON serialization to keep
witnessed outputs stable across runs.

\section{Threat model}
Cogitator detects accidental drift and supports forensic diagnosis.
It does not provide adversarial tamper resistance unless expected roots are
anchored in a trusted channel (e.g., signed releases or an external transparency log).
An attacker who can modify both artifacts and expected roots can bypass detection.

\section{System design}
A run is defined by inputs (task set, seed, flags), plus optional agent mode
configuration. The harness emits a run bundle of artifacts and two roots
(Figure~\ref{fig:pipeline}).

\paragraph{Determinism caveat.}
Semantic determinism is guaranteed for the same binary and the same inputs.
Cross-platform bit-for-bit identity may require additional constraints
(e.g., consistent floating-point behavior and locale settings).

\begin{figure}[t]
  \centering
  \begin{tikzpicture}[font=\small, node distance=1.6cm, >=latex]
  \tikzstyle{block}=[draw, rounded corners, align=center, minimum height=1cm, minimum width=2.8cm]

  \node[block] (inputs) {Inputs\\(seed, flags, tasks)};
  \node[block, right=of inputs] (run) {Deterministic\\run};
  \node[block, right=of run] (trace) {Trace / transcript};
  \node[block, below=of trace] (artifacts) {Artifacts\\(bundle)};
  \node[block, right=of trace] (roots) {Roots\\\ExecutionRoot\\\BundleRoot};
  \node[block, right=of roots] (verify) {Verify\\(CI + drift)};

  \draw[->] (inputs) -- (run);
  \draw[->] (run) -- (trace);
  \draw[->] (trace) -- (roots);
  \draw[->] (trace) -- (artifacts);
  \draw[->] (artifacts) -- (roots);
  \draw[->] (roots) -- (verify);
\end{tikzpicture}

  \caption{Cogitator pipeline: inputs drive a run that emits a trace/transcript and artifacts.
  \ExecutionRoot binds to witnessed metadata and semantic events, while \BundleRoot binds to the
  full artifact bundle for integrity verification.}
  \label{fig:pipeline}
\end{figure}

\section{Witness model}
Cogitator distinguishes \emph{witnessed} metadata (inputs to the root) from
\emph{provenance} metadata (environment context). Only witnessed metadata influences
\ExecutionRoot. Provenance is retained for auditability but excluded from root
computation to preserve cross-machine stability.

\subsection{Witnessed metadata}
Witnessed metadata is serialized canonically and includes:
\begin{itemize}[leftmargin=1.4em]
  \item \texttt{schema\_version} (trace schema)
  \item \texttt{seed}, \texttt{requested\_runs}, \texttt{executed\_runs}
  \item \texttt{parallel} and \texttt{parallel\_strategy}
  \item \texttt{case\_filter}
  \item \texttt{entropy\_sources} and \texttt{total\_rng\_calls}
  \item \texttt{chaos\_profile} summary (when fault injection is enabled)
  \item \texttt{pass\_threshold} (when a gauntlet threshold is configured)
\end{itemize}

\subsection{Provenance metadata}
Provenance metadata is recorded for debugging and audit trails and includes:
\begin{itemize}[leftmargin=1.4em]
  \item \texttt{created\_at}
  \item \texttt{git\_rev}, \texttt{rustc\_version}, \texttt{cargo\_version}
  \item \texttt{nix\_store\_path} and optional \texttt{nix\_provenance}
  \item \texttt{agent\_threads} (agent mode throughput setting)
  \item \texttt{variability\_factors}
\end{itemize}

\section{Cryptographic commitments}
Cogitator uses two roots with distinct purposes. \ExecutionRoot captures semantic
execution; \BundleRoot captures artifact integrity.

\subsection{ExecutionRoot}
\ExecutionRoot commits to witnessed metadata and a deterministic semantic event stream.
The event stream is serialized in canonical JSON and fed in order to a BLAKE3-based
hash accumulator.

\paragraph{Non-agent mode.}
\ExecutionRoot is computed from witnessed metadata plus all trace events ordered by
\texttt{(run\_id, step)}. Each trace event includes the schema version, run identifier,
case identifier, step number, role, content, and entropy counters.

\paragraph{Agent mode.}
\ExecutionRoot is computed from witnessed metadata, then the agent trace in order, and
then tool transcript entries per step ordered by \texttt{tool\_call\_idx}. Tool call
witness entries include the step, tool name, request, outcome, and fault annotation.

\subsection{BundleRoot}
\BundleRoot commits to the artifact bundle. The harness constructs a witness manifest
that records per-artifact BLAKE3 hashes and the canonical bundle hash. The bundle hash
is computed as the BLAKE3 hash of the canonical JSON encoding of the
\texttt{artifact\_hashes} map, and it is stored as \texttt{bundle\_hash} in the manifest.

\section{Artifacts and modes}
Table~\ref{tab:artifacts} summarizes the core (non-agent) and agent-mode artifacts.
All artifacts use canonical JSON encoding where applicable.

\begin{table}[t]
  \centering
  \caption{Run artifacts emitted in core and agent modes. Optional artifacts are marked
  with \textsuperscript{*}.}
  \label{tab:artifacts}
  \begin{tabular}{@{}p{0.28\linewidth}p{0.32\linewidth}p{0.32\linewidth}@{}}
    \toprule
    Artifact & Core mode & Agent mode \\
    \midrule
    \texttt{meta.json} & witnessed + provenance & witnessed + provenance \\
    \texttt{trace.jsonl} & trace events & --- \\
    \texttt{results.csv/json} & case results & --- \\
    \texttt{summary.json} & aggregate metrics & --- \\
    \texttt{analysis.json} & metadata + summary + results & --- \\
    \texttt{witness\_root.txt} & \ExecutionRoot & \ExecutionRoot for agent trace \\
    \texttt{agent\_trace.json} & --- & ordered agent trace \\
    \texttt{tool\_transcript.json} & --- & tool calls and outcomes \\
    \texttt{hash\_chain.txt} & --- & linear hash chain \\
    \texttt{drift\_report.json} & --- & drift report \\
    \texttt{chaos\_profile.json}\textsuperscript{*} & --- & fault schedule \\
    \texttt{nix\_provenance.json}\textsuperscript{*} & optional & optional \\
    \texttt{witness\_manifest.json} & --- & manifest + \BundleRoot \\
    \bottomrule
  \end{tabular}
\end{table}

\section{Agent replay and drift}
Agent mode treats tool calls as a record/replay boundary. In live runs, tool responses
are captured in the transcript. In replay runs, responses are replayed verbatim, making
agent traces deterministic for a fixed seed and transcript. Drift detection compares
transcripts, tool call ordering, and the resulting \ExecutionRoot to surface semantic
changes.

\section{Fault injection}
Cogitator supports deterministic fault injection at tool boundaries. A seeded fault
schedule decides timeouts, drops, corruptions, or latency simulations and is recorded
in \texttt{chaos\_profile.json}. The chaos summary is included in witnessed metadata,
while simulated latency itself is excluded from witness commitments by default.

\section{Verification}
Verification has two complementary checks:
\begin{enumerate}[leftmargin=1.4em]
  \item \textbf{Bundle integrity:} recompute per-artifact hashes, verify them against the
  manifest, and recompute \BundleRoot from the canonical JSON \texttt{artifact\_hashes}.
  This is the comparison used for CI gating.
  \item \textbf{Execution integrity:} recompute \ExecutionRoot from witnessed metadata and
  the semantic event stream (trace for core runs, agent trace + tool transcript for
  agent runs). This detects semantic drift even when the bundle is intact.
\end{enumerate}
CI should compare \BundleRoot for artifact integrity, while \ExecutionRoot highlights
semantic drift in the witnessed execution stream.

\section{Evaluation}
Cogitator includes unit tests and property-based tests for determinism invariants.
At scale, the harness has been exercised on a large run (e.g., 500{,}000 cases) with
stable \ExecutionRoot and \BundleRoot across repeated runs under identical inputs.
Dataset fingerprint (SHA-256): \texttt{TODO}.

\section{Related work}
Cogitator aligns with prior work on deterministic simulation testing, reproducibility
metadata, and canonical JSON serialization.\cite{TODO-dst,TODO-repro,TODO-jcs}

\section{Limitations and future work}
Future work includes Merkle-tree inclusion proofs for artifact subsets, cross-platform
float normalization policies, and signed attestations for expected roots.

\appendix
\section{Worked example}
Consider a core run with witnessed metadata $M$ and trace events $E$ ordered by
\texttt{(run\_id, step)}. \ExecutionRoot is computed by canonicalizing $M$, initializing
an accumulator, and then hashing each event in order. For an agent run, the event stream
is replaced by the agent trace followed by tool calls ordered by \texttt{tool\_call\_idx}.

For \BundleRoot, the witness manifest records artifact hashes such as
\texttt{meta.json} $\rightarrow h_1$, \texttt{trace.jsonl} $\rightarrow h_2$, and so on.
The bundle hash is computed as the hash of the canonical JSON map
\texttt{\{artifact\_hashes\}} and stored in \texttt{witness\_manifest.json}. Verification
recomputes the artifact hashes and bundle hash and compares them to the manifest.

\bibliographystyle{plain}
\bibliography{references}

\end{document}
